\section{Conclusion}
\label{sec:conclusion}

The project set out to reproduce a published amateur-radio OFDM PHY
layer and ended with a complete adaptive communication system.  The
reproduction itself is closed: 44/44 worked examples bit-exact, the
BER/PER figures reproduced, and the sensitivity gap traced to three
specific synchronization defects and eliminated to within 0.3\,dB of a
genie-synchronized bound.  On that foundation, the waveform was
extended an order of magnitude in each direction --- down to
$-17.9$\,dB and 7.8\,bit/s (78\,\% of Shannon capacity at $-20$\,dB)
and up to $+4.7$\,dB and 1059\,bit/s with 16-QAM --- and wrapped in a
link layer that measures everything it can observe and feeds all of it
back: receiver-driven rate adaptation over a 13-rung measured ladder,
CRC-gated HARQ from exported LLRs, learned per-rung sensitivity
offsets, and closed-loop AFC netting of the stations' carrier
frequencies.

Two results deserve emphasis beyond the feature list.  First, the
measured LLR recalibration --- 1.5--2\,dB from correcting the
\emph{shape} of the demodulator's confidence, at zero over-the-air
cost --- illustrates how much sensitivity can hide between a correct
demodulator and a correct decoder.  Second, the fixed-point receiver
matching and then beating the floating-point chain at the sensitivity
edge shows that an RTL-oriented implementation need not be a
degradation; the discipline it forces (bounded scales, division-free
estimation, a tracker instead of a polynomial fit) produced the better
algorithm.

Everything claimed here is regenerable from the repository: the
regression suites are self-verifying, every sensitivity figure has its
sweep script and JSON dump, and the final demonstration is a WAV file
any copy of the receiver decodes blind.  The natural next step is the
one simulation cannot take: the same frames over a real HF path.
